\subsection{Exploratory sensitivity to the numerical lattice}

Exploratory two-dimensional calculations at the seventh radial location
also indicate sensitivity to numerical spatial discretization. Here the
width is the amplitude decay rate $\Gamma=-\operatorname{Im}\rho$ for
the convention $\exp(-i\rho t)$. The reported minimum widths scale
approximately as $h^4$ for a five-point square stencil and $h^8$ for both
a weighted nine-point square stencil and a seven-point triangular stencil.
The leading angular radiation components are $\cos(4\theta)$ for both
square stencils and $\cos(6\theta)$ for the triangular stencil. These
observations are consistent with a perturbative interpretation in which
the leading anisotropic operator correction induces a radiation amplitude
linear in its strength, and the width is quadratic in that amplitude.
The analogous quadratic asymmetry law is established for optical
quasi-BICs by Koshelev et al.~\cite{Koshelev2018}, Eqs.~(2)--(3);
its application here remains a model-specific interpretation.
The reported exponent uncertainties are fit errors, and comparisons of
absorbing layers and shifted eigenvalue solves are sensitivity diagnostics
rather than certified error bounds. The calculations concern one
two-dimensional mode and do not establish three-dimensional or nonlinear
persistence. They motivate explicit stencil and resolution checks before
angular radiation patterns are interpreted as physical structure.
